% GENERATED BY src/python/lmi/tables/ext_asyl_threeway.py
% -- DO NOT EDIT
%==============================================================================%
% ext_asyl_threeway.tex
%
% Purpose : Erweiterung, Tabelle E4 (Anhang) -- explorative Dreifach-
%           interaktion Kursangebot x Aufenthaltsdauer x Asylstatus.
% Source  : SOEPremote job 244839, src/stata/remote/ext_asyl_threeway.do
%           Rohlisting: results/transcripts/
%                       src_stata_remote_ext_asyl_threeway.do.eml
%           Transkript: results/comparisons/ext_asyl_employment.md (Anhang)
% Needs   : booktabs, tabularx (beide bereits in bachelorarbeit.sty geladen)
% Usage   : % GENERATED BY src/python/lmi/tables/ext_asyl_threeway.py
% -- DO NOT EDIT
%==============================================================================%
% ext_asyl_threeway.tex
%
% Purpose : Erweiterung, Tabelle E4 (Anhang) -- explorative Dreifach-
%           interaktion Kursangebot x Aufenthaltsdauer x Asylstatus.
% Source  : SOEPremote job 244839, src/stata/remote/ext_asyl_threeway.do
%           Rohlisting: results/transcripts/
%                       src_stata_remote_ext_asyl_threeway.do.eml
%           Transkript: results/comparisons/ext_asyl_employment.md (Anhang)
% Needs   : booktabs, tabularx (beide bereits in bachelorarbeit.sty geladen)
% Usage   : % GENERATED BY src/python/lmi/tables/ext_asyl_threeway.py
% -- DO NOT EDIT
%==============================================================================%
% ext_asyl_threeway.tex
%
% Purpose : Erweiterung, Tabelle E4 (Anhang) -- explorative Dreifach-
%           interaktion Kursangebot x Aufenthaltsdauer x Asylstatus.
% Source  : SOEPremote job 244839, src/stata/remote/ext_asyl_threeway.do
%           Rohlisting: results/transcripts/
%                       src_stata_remote_ext_asyl_threeway.do.eml
%           Transkript: results/comparisons/ext_asyl_employment.md (Anhang)
% Needs   : booktabs, tabularx (beide bereits in bachelorarbeit.sty geladen)
% Usage   : % GENERATED BY src/python/lmi/tables/ext_asyl_threeway.py
% -- DO NOT EDIT
%==============================================================================%
% ext_asyl_threeway.tex
%
% Purpose : Erweiterung, Tabelle E4 (Anhang) -- explorative Dreifach-
%           interaktion Kursangebot x Aufenthaltsdauer x Asylstatus.
% Source  : SOEPremote job 244839, src/stata/remote/ext_asyl_threeway.do
%           Rohlisting: results/transcripts/
%                       src_stata_remote_ext_asyl_threeway.do.eml
%           Transkript: results/comparisons/ext_asyl_employment.md (Anhang)
% Needs   : booktabs, tabularx (beide bereits in bachelorarbeit.sty geladen)
% Usage   : \input{tables/ext_asyl_threeway}
% Author  : Emre Suemer, 2026-09-01
%==============================================================================%

\makeatletter
\@ifundefined{NC@rewrite@Z}{%
  \newcolumntype{Z}{>{\raggedright\arraybackslash\hangindent=1.9em\hangafter=1}X}%
}{}
\makeatother

\begin{table}[htbp]
  \centering
  \caption{Explorativ: Kursangebot, Aufenthaltsdauer und Asylstatus}
  \label{tab:ext-asyl-dreifach}
  \footnotesize
  \setlength{\tabcolsep}{4pt}%
  \begin{tabularx}{\textwidth}{@{}Zrr@{}}
    \toprule
     & Modell E4.1 & Modell E4.2 \\
     & Erwerbst\"atigkeit & Deutschkenntnisse \\
    \midrule
    \textbf{Kreisangebot an Integrationskursen, std.} & 0{,}0028 & 0{,}1493 \\
     & (0{,}0225) & (0{,}1471) \\[0.3ex]
    Aufenthaltsdauer, in Monaten (/12) & 0{,}1016$^{**}$ & 0{,}9943$^{**}$ \\
     & (0{,}0181) & (0{,}1123) \\[0.3ex]
    \quad Angebot $\times$ Aufenthaltsdauer & 0{,}0164 & $-$0{,}0443 \\
     & (0{,}0121) & (0{,}0703) \\[0.3ex]
    \quad Angebot $\times$ anerkannt & $-$0{,}0326 & 0{,}3114 \\
     & (0{,}0372) & (0{,}2490) \\[0.3ex]
    \quad Angebot $\times$ abgelehnt & 0{,}0407 & 0{,}0913 \\
     & (0{,}0440) & (0{,}2973) \\[0.3ex]
    \quad Aufenthaltsdauer $\times$ anerkannt & 0{,}0129 & $-$0{,}0435 \\
     & (0{,}0169) & (0{,}0952) \\[0.3ex]
    \quad Aufenthaltsdauer $\times$ abgelehnt & 0{,}0181 & $-$0{,}1120 \\
     & (0{,}0188) & (0{,}1116) \\[0.3ex]
    \quad Angebot $\times$ Aufenthaltsdauer $\times$ anerkannt & 0{,}0040 & $-$0{,}1361 \\
     & (0{,}0161) & (0{,}0958) \\[0.3ex]
    \quad Angebot $\times$ Aufenthaltsdauer $\times$ abgelehnt & $-$0{,}0265 & $-$0{,}0431 \\
     & (0{,}0192) & (0{,}1079) \\[0.3ex]
    Asylstatus: anerkannt & 0{,}0195 & 0{,}9488$^{**}$ \\
     & (0{,}0429) & (0{,}2716) \\[0.3ex]
    Asylstatus: abgelehnt & 0{,}0333 & 0{,}1581 \\
     & (0{,}0472) & (0{,}3278) \\[0.3ex]
    Konstante & $-$0{,}0196 & 4{,}7547$^{**}$ \\
     & (0{,}0639) & (0{,}4213) \\
    \midrule
    Gemeinsamer Test der beiden Dreifachterme & F(2; 2727{,}0)\,=\,1{,}40 & F(2; 2711{,}8)\,=\,1{,}05 \\
    \quad p & 0{,}247 & 0{,}350 \\
    \midrule
    N Beobachtungen & 5.467 & 5.467 \\
    N Personen & 2.730 & 2.730 \\
    N Imputationen & 20 & 20 \\
    R\textsuperscript{2} (adj.) & 0{,}2152 & 0{,}3912 \\
    \bottomrule
  \end{tabularx}

  \vspace{0.75ex}
  \parbox{\textwidth}{\footnotesize
    \textit{Anmerkungen:} Eigene Berechnung, IAB-BAMF-SOEP-Befragung von
    Gefl\"uchteten 2016--2019, SOEP-Core v36 (Remote Edition), gewichtet.
    Spezifikation wie Modell~1.1 in Tabelle~\ref{tab:erwerbstaetigkeit},
    erg\"anzt um die vollst\"andige Dreifachinteraktion
    \texttt{ib1.asyl\_req\_stat\#\#c.z\_ratio\_courseoport\#\#c.stay\_dur}.
	%    Individuelle und regionale Kontrollvariablen sowie fixe Effekte f\"ur
	%    Erhebungsjahr, Erhebungsstichprobe, Bundesland und Herkunftsland sind
	%    enthalten und nicht ausgewiesen; die vollst\"andigen Koeffizienten stehen
	%    in \texttt{build/tables/ext\_asyl\_threeway.csv}.
    
    \textbf{Rohe Einheiten, keine Prozentpunkte} -- anders als
    Tabelle~\ref{tab:ext-asyl-erwerb}, die mit 100 multipliziert ist. Die
    Spalte Erwerbst\"atigkeit ist deshalb nicht zellenweise mit ihr
    vergleichbar (0{,}01\,=\,1~p.p.). Robuste, auf Personenebene geclusterte
    Standardfehler in Klammern.
    Referenzkategorie des Asylstatus ist \glqq keine Entscheidung\grqq{}.
    $^{+}$~p\,$<$\,0{,}1, $^{*}$~p\,$<$\,0{,}05, $^{**}$~p\,$<$\,0{,}01
    (zweiseitig).
%    \\[0.4ex]
%    \textbf{Explorativ, und vorab als solches festgelegt.} Getestet wird, ob
%    der in der Vorlage berichtete R\"uckgang des Angebotseffekts mit
%    zunehmender Aufenthaltsdauer dort konzentriert ist, wo der Zeithorizont
%    fr\"uh durch eine Entscheidung fixiert wird. Das Designdokument hielt vor
%    dem Lauf fest, dass diese Spezifikation mit hoher Wahrscheinlichkeit
%    unterpowert ist; genau das zeigt sich: beide gemeinsamen Tests sind
%    insignifikant (p\,=\,0{,}247 und p\,=\,0{,}350), und jeder Standardfehler
%    der Dreifachterme ist gro\ss{} gegen\"uber dem zugeh\"origen
%    Koeffizienten. Aus dieser Tabelle sind ausschlie\ss{}lich die beiden
%    gemeinsamen p-Werte zu berichten; einzelne Koeffizienten tragen keine
%    Interpretation.
    
    }
\end{table}

% Author  : Emre Suemer, 2026-09-01
%==============================================================================%

\makeatletter
\@ifundefined{NC@rewrite@Z}{%
  \newcolumntype{Z}{>{\raggedright\arraybackslash\hangindent=1.9em\hangafter=1}X}%
}{}
\makeatother

\begin{table}[htbp]
  \centering
  \caption{Explorativ: Kursangebot, Aufenthaltsdauer und Asylstatus}
  \label{tab:ext-asyl-dreifach}
  \footnotesize
  \setlength{\tabcolsep}{4pt}%
  \begin{tabularx}{\textwidth}{@{}Zrr@{}}
    \toprule
     & Modell E4.1 & Modell E4.2 \\
     & Erwerbst\"atigkeit & Deutschkenntnisse \\
    \midrule
    \textbf{Kreisangebot an Integrationskursen, std.} & 0{,}0028 & 0{,}1493 \\
     & (0{,}0225) & (0{,}1471) \\[0.3ex]
    Aufenthaltsdauer, in Monaten (/12) & 0{,}1016$^{**}$ & 0{,}9943$^{**}$ \\
     & (0{,}0181) & (0{,}1123) \\[0.3ex]
    \quad Angebot $\times$ Aufenthaltsdauer & 0{,}0164 & $-$0{,}0443 \\
     & (0{,}0121) & (0{,}0703) \\[0.3ex]
    \quad Angebot $\times$ anerkannt & $-$0{,}0326 & 0{,}3114 \\
     & (0{,}0372) & (0{,}2490) \\[0.3ex]
    \quad Angebot $\times$ abgelehnt & 0{,}0407 & 0{,}0913 \\
     & (0{,}0440) & (0{,}2973) \\[0.3ex]
    \quad Aufenthaltsdauer $\times$ anerkannt & 0{,}0129 & $-$0{,}0435 \\
     & (0{,}0169) & (0{,}0952) \\[0.3ex]
    \quad Aufenthaltsdauer $\times$ abgelehnt & 0{,}0181 & $-$0{,}1120 \\
     & (0{,}0188) & (0{,}1116) \\[0.3ex]
    \quad Angebot $\times$ Aufenthaltsdauer $\times$ anerkannt & 0{,}0040 & $-$0{,}1361 \\
     & (0{,}0161) & (0{,}0958) \\[0.3ex]
    \quad Angebot $\times$ Aufenthaltsdauer $\times$ abgelehnt & $-$0{,}0265 & $-$0{,}0431 \\
     & (0{,}0192) & (0{,}1079) \\[0.3ex]
    Asylstatus: anerkannt & 0{,}0195 & 0{,}9488$^{**}$ \\
     & (0{,}0429) & (0{,}2716) \\[0.3ex]
    Asylstatus: abgelehnt & 0{,}0333 & 0{,}1581 \\
     & (0{,}0472) & (0{,}3278) \\[0.3ex]
    Konstante & $-$0{,}0196 & 4{,}7547$^{**}$ \\
     & (0{,}0639) & (0{,}4213) \\
    \midrule
    Gemeinsamer Test der beiden Dreifachterme & F(2; 2727{,}0)\,=\,1{,}40 & F(2; 2711{,}8)\,=\,1{,}05 \\
    \quad p & 0{,}247 & 0{,}350 \\
    \midrule
    N Beobachtungen & 5.467 & 5.467 \\
    N Personen & 2.730 & 2.730 \\
    N Imputationen & 20 & 20 \\
    R\textsuperscript{2} (adj.) & 0{,}2152 & 0{,}3912 \\
    \bottomrule
  \end{tabularx}

  \vspace{0.75ex}
  \parbox{\textwidth}{\footnotesize
    \textit{Anmerkungen:} Eigene Berechnung, IAB-BAMF-SOEP-Befragung von
    Gefl\"uchteten 2016--2019, SOEP-Core v36 (Remote Edition), gewichtet.
    Spezifikation wie Modell~1.1 in Tabelle~\ref{tab:erwerbstaetigkeit},
    erg\"anzt um die vollst\"andige Dreifachinteraktion
    \texttt{ib1.asyl\_req\_stat\#\#c.z\_ratio\_courseoport\#\#c.stay\_dur}.
	%    Individuelle und regionale Kontrollvariablen sowie fixe Effekte f\"ur
	%    Erhebungsjahr, Erhebungsstichprobe, Bundesland und Herkunftsland sind
	%    enthalten und nicht ausgewiesen; die vollst\"andigen Koeffizienten stehen
	%    in \texttt{build/tables/ext\_asyl\_threeway.csv}.
    
    \textbf{Rohe Einheiten, keine Prozentpunkte} -- anders als
    Tabelle~\ref{tab:ext-asyl-erwerb}, die mit 100 multipliziert ist. Die
    Spalte Erwerbst\"atigkeit ist deshalb nicht zellenweise mit ihr
    vergleichbar (0{,}01\,=\,1~p.p.). Robuste, auf Personenebene geclusterte
    Standardfehler in Klammern.
    Referenzkategorie des Asylstatus ist \glqq keine Entscheidung\grqq{}.
    $^{+}$~p\,$<$\,0{,}1, $^{*}$~p\,$<$\,0{,}05, $^{**}$~p\,$<$\,0{,}01
    (zweiseitig).
%    \\[0.4ex]
%    \textbf{Explorativ, und vorab als solches festgelegt.} Getestet wird, ob
%    der in der Vorlage berichtete R\"uckgang des Angebotseffekts mit
%    zunehmender Aufenthaltsdauer dort konzentriert ist, wo der Zeithorizont
%    fr\"uh durch eine Entscheidung fixiert wird. Das Designdokument hielt vor
%    dem Lauf fest, dass diese Spezifikation mit hoher Wahrscheinlichkeit
%    unterpowert ist; genau das zeigt sich: beide gemeinsamen Tests sind
%    insignifikant (p\,=\,0{,}247 und p\,=\,0{,}350), und jeder Standardfehler
%    der Dreifachterme ist gro\ss{} gegen\"uber dem zugeh\"origen
%    Koeffizienten. Aus dieser Tabelle sind ausschlie\ss{}lich die beiden
%    gemeinsamen p-Werte zu berichten; einzelne Koeffizienten tragen keine
%    Interpretation.
    
    }
\end{table}

% Author  : Emre Suemer, 2026-09-01
%==============================================================================%

\makeatletter
\@ifundefined{NC@rewrite@Z}{%
  \newcolumntype{Z}{>{\raggedright\arraybackslash\hangindent=1.9em\hangafter=1}X}%
}{}
\makeatother

\begin{table}[htbp]
  \centering
  \caption{Explorativ: Kursangebot, Aufenthaltsdauer und Asylstatus}
  \label{tab:ext-asyl-dreifach}
  \footnotesize
  \setlength{\tabcolsep}{4pt}%
  \begin{tabularx}{\textwidth}{@{}Zrr@{}}
    \toprule
     & Modell E4.1 & Modell E4.2 \\
     & Erwerbst\"atigkeit & Deutschkenntnisse \\
    \midrule
    \textbf{Kreisangebot an Integrationskursen, std.} & 0{,}0028 & 0{,}1493 \\
     & (0{,}0225) & (0{,}1471) \\[0.3ex]
    Aufenthaltsdauer, in Monaten (/12) & 0{,}1016$^{**}$ & 0{,}9943$^{**}$ \\
     & (0{,}0181) & (0{,}1123) \\[0.3ex]
    \quad Angebot $\times$ Aufenthaltsdauer & 0{,}0164 & $-$0{,}0443 \\
     & (0{,}0121) & (0{,}0703) \\[0.3ex]
    \quad Angebot $\times$ anerkannt & $-$0{,}0326 & 0{,}3114 \\
     & (0{,}0372) & (0{,}2490) \\[0.3ex]
    \quad Angebot $\times$ abgelehnt & 0{,}0407 & 0{,}0913 \\
     & (0{,}0440) & (0{,}2973) \\[0.3ex]
    \quad Aufenthaltsdauer $\times$ anerkannt & 0{,}0129 & $-$0{,}0435 \\
     & (0{,}0169) & (0{,}0952) \\[0.3ex]
    \quad Aufenthaltsdauer $\times$ abgelehnt & 0{,}0181 & $-$0{,}1120 \\
     & (0{,}0188) & (0{,}1116) \\[0.3ex]
    \quad Angebot $\times$ Aufenthaltsdauer $\times$ anerkannt & 0{,}0040 & $-$0{,}1361 \\
     & (0{,}0161) & (0{,}0958) \\[0.3ex]
    \quad Angebot $\times$ Aufenthaltsdauer $\times$ abgelehnt & $-$0{,}0265 & $-$0{,}0431 \\
     & (0{,}0192) & (0{,}1079) \\[0.3ex]
    Asylstatus: anerkannt & 0{,}0195 & 0{,}9488$^{**}$ \\
     & (0{,}0429) & (0{,}2716) \\[0.3ex]
    Asylstatus: abgelehnt & 0{,}0333 & 0{,}1581 \\
     & (0{,}0472) & (0{,}3278) \\[0.3ex]
    Konstante & $-$0{,}0196 & 4{,}7547$^{**}$ \\
     & (0{,}0639) & (0{,}4213) \\
    \midrule
    Gemeinsamer Test der beiden Dreifachterme & F(2; 2727{,}0)\,=\,1{,}40 & F(2; 2711{,}8)\,=\,1{,}05 \\
    \quad p & 0{,}247 & 0{,}350 \\
    \midrule
    N Beobachtungen & 5.467 & 5.467 \\
    N Personen & 2.730 & 2.730 \\
    N Imputationen & 20 & 20 \\
    R\textsuperscript{2} (adj.) & 0{,}2152 & 0{,}3912 \\
    \bottomrule
  \end{tabularx}

  \vspace{0.75ex}
  \parbox{\textwidth}{\footnotesize
    \textit{Anmerkungen:} Eigene Berechnung, IAB-BAMF-SOEP-Befragung von
    Gefl\"uchteten 2016--2019, SOEP-Core v36 (Remote Edition), gewichtet.
    Spezifikation wie Modell~1.1 in Tabelle~\ref{tab:erwerbstaetigkeit},
    erg\"anzt um die vollst\"andige Dreifachinteraktion
    \texttt{ib1.asyl\_req\_stat\#\#c.z\_ratio\_courseoport\#\#c.stay\_dur}.
	%    Individuelle und regionale Kontrollvariablen sowie fixe Effekte f\"ur
	%    Erhebungsjahr, Erhebungsstichprobe, Bundesland und Herkunftsland sind
	%    enthalten und nicht ausgewiesen; die vollst\"andigen Koeffizienten stehen
	%    in \texttt{build/tables/ext\_asyl\_threeway.csv}.
    
    \textbf{Rohe Einheiten, keine Prozentpunkte} -- anders als
    Tabelle~\ref{tab:ext-asyl-erwerb}, die mit 100 multipliziert ist. Die
    Spalte Erwerbst\"atigkeit ist deshalb nicht zellenweise mit ihr
    vergleichbar (0{,}01\,=\,1~p.p.). Robuste, auf Personenebene geclusterte
    Standardfehler in Klammern.
    Referenzkategorie des Asylstatus ist \glqq keine Entscheidung\grqq{}.
    $^{+}$~p\,$<$\,0{,}1, $^{*}$~p\,$<$\,0{,}05, $^{**}$~p\,$<$\,0{,}01
    (zweiseitig).
%    \\[0.4ex]
%    \textbf{Explorativ, und vorab als solches festgelegt.} Getestet wird, ob
%    der in der Vorlage berichtete R\"uckgang des Angebotseffekts mit
%    zunehmender Aufenthaltsdauer dort konzentriert ist, wo der Zeithorizont
%    fr\"uh durch eine Entscheidung fixiert wird. Das Designdokument hielt vor
%    dem Lauf fest, dass diese Spezifikation mit hoher Wahrscheinlichkeit
%    unterpowert ist; genau das zeigt sich: beide gemeinsamen Tests sind
%    insignifikant (p\,=\,0{,}247 und p\,=\,0{,}350), und jeder Standardfehler
%    der Dreifachterme ist gro\ss{} gegen\"uber dem zugeh\"origen
%    Koeffizienten. Aus dieser Tabelle sind ausschlie\ss{}lich die beiden
%    gemeinsamen p-Werte zu berichten; einzelne Koeffizienten tragen keine
%    Interpretation.
    
    }
\end{table}

% Author  : Emre Suemer, 2026-09-01
%==============================================================================%

\makeatletter
\@ifundefined{NC@rewrite@Z}{%
  \newcolumntype{Z}{>{\raggedright\arraybackslash\hangindent=1.9em\hangafter=1}X}%
}{}
\makeatother

\begin{table}[htbp]
  \centering
  \caption{Explorativ: Kursangebot, Aufenthaltsdauer und Asylstatus}
  \label{tab:ext-asyl-dreifach}
  \footnotesize
  \setlength{\tabcolsep}{4pt}%
  \begin{tabularx}{\textwidth}{@{}Zrr@{}}
    \toprule
     & Modell E4.1 & Modell E4.2 \\
     & Erwerbst\"atigkeit & Deutschkenntnisse \\
    \midrule
    \textbf{Kreisangebot an Integrationskursen, std.} & 0{,}0028 & 0{,}1493 \\
     & (0{,}0225) & (0{,}1471) \\[0.3ex]
    Aufenthaltsdauer, in Monaten (/12) & 0{,}1016$^{**}$ & 0{,}9943$^{**}$ \\
     & (0{,}0181) & (0{,}1123) \\[0.3ex]
    \quad Angebot $\times$ Aufenthaltsdauer & 0{,}0164 & $-$0{,}0443 \\
     & (0{,}0121) & (0{,}0703) \\[0.3ex]
    \quad Angebot $\times$ anerkannt & $-$0{,}0326 & 0{,}3114 \\
     & (0{,}0372) & (0{,}2490) \\[0.3ex]
    \quad Angebot $\times$ abgelehnt & 0{,}0407 & 0{,}0913 \\
     & (0{,}0440) & (0{,}2973) \\[0.3ex]
    \quad Aufenthaltsdauer $\times$ anerkannt & 0{,}0129 & $-$0{,}0435 \\
     & (0{,}0169) & (0{,}0952) \\[0.3ex]
    \quad Aufenthaltsdauer $\times$ abgelehnt & 0{,}0181 & $-$0{,}1120 \\
     & (0{,}0188) & (0{,}1116) \\[0.3ex]
    \quad Angebot $\times$ Aufenthaltsdauer $\times$ anerkannt & 0{,}0040 & $-$0{,}1361 \\
     & (0{,}0161) & (0{,}0958) \\[0.3ex]
    \quad Angebot $\times$ Aufenthaltsdauer $\times$ abgelehnt & $-$0{,}0265 & $-$0{,}0431 \\
     & (0{,}0192) & (0{,}1079) \\[0.3ex]
    Asylstatus: anerkannt & 0{,}0195 & 0{,}9488$^{**}$ \\
     & (0{,}0429) & (0{,}2716) \\[0.3ex]
    Asylstatus: abgelehnt & 0{,}0333 & 0{,}1581 \\
     & (0{,}0472) & (0{,}3278) \\[0.3ex]
    Konstante & $-$0{,}0196 & 4{,}7547$^{**}$ \\
     & (0{,}0639) & (0{,}4213) \\
    \midrule
    Gemeinsamer Test der beiden Dreifachterme & F(2; 2727{,}0)\,=\,1{,}40 & F(2; 2711{,}8)\,=\,1{,}05 \\
    \quad p & 0{,}247 & 0{,}350 \\
    \midrule
    N Beobachtungen & 5.467 & 5.467 \\
    N Personen & 2.730 & 2.730 \\
    N Imputationen & 20 & 20 \\
    R\textsuperscript{2} (adj.) & 0{,}2152 & 0{,}3912 \\
    \bottomrule
  \end{tabularx}

  \vspace{0.75ex}
  \parbox{\textwidth}{\footnotesize
    \textit{Anmerkungen:} Eigene Berechnung, IAB-BAMF-SOEP-Befragung von
    Gefl\"uchteten 2016--2019, SOEP-Core v36 (Remote Edition), gewichtet.
    Spezifikation wie Modell~1.1 in Tabelle~\ref{tab:erwerbstaetigkeit},
    erg\"anzt um die vollst\"andige Dreifachinteraktion
    \texttt{ib1.asyl\_req\_stat\#\#c.z\_ratio\_courseoport\#\#c.stay\_dur}.
	%    Individuelle und regionale Kontrollvariablen sowie fixe Effekte f\"ur
	%    Erhebungsjahr, Erhebungsstichprobe, Bundesland und Herkunftsland sind
	%    enthalten und nicht ausgewiesen; die vollst\"andigen Koeffizienten stehen
	%    in \texttt{build/tables/ext\_asyl\_threeway.csv}.
    
    \textbf{Rohe Einheiten, keine Prozentpunkte} -- anders als
    Tabelle~\ref{tab:ext-asyl-erwerb}, die mit 100 multipliziert ist. Die
    Spalte Erwerbst\"atigkeit ist deshalb nicht zellenweise mit ihr
    vergleichbar (0{,}01\,=\,1~p.p.). Robuste, auf Personenebene geclusterte
    Standardfehler in Klammern.
    Referenzkategorie des Asylstatus ist \glqq keine Entscheidung\grqq{}.
    $^{+}$~p\,$<$\,0{,}1, $^{*}$~p\,$<$\,0{,}05, $^{**}$~p\,$<$\,0{,}01
    (zweiseitig).
%    \\[0.4ex]
%    \textbf{Explorativ, und vorab als solches festgelegt.} Getestet wird, ob
%    der in der Vorlage berichtete R\"uckgang des Angebotseffekts mit
%    zunehmender Aufenthaltsdauer dort konzentriert ist, wo der Zeithorizont
%    fr\"uh durch eine Entscheidung fixiert wird. Das Designdokument hielt vor
%    dem Lauf fest, dass diese Spezifikation mit hoher Wahrscheinlichkeit
%    unterpowert ist; genau das zeigt sich: beide gemeinsamen Tests sind
%    insignifikant (p\,=\,0{,}247 und p\,=\,0{,}350), und jeder Standardfehler
%    der Dreifachterme ist gro\ss{} gegen\"uber dem zugeh\"origen
%    Koeffizienten. Aus dieser Tabelle sind ausschlie\ss{}lich die beiden
%    gemeinsamen p-Werte zu berichten; einzelne Koeffizienten tragen keine
%    Interpretation.
    
    }
\end{table}
